% !TEX root = ../main.tex

\section{Model v1.1: A Deep-Supervised Multi-Decoder CNN}
\label{sec:model11}

Model v1.1 is the fully-convolutional model resulting from the iterations in \secref{subsec:beta-v10}--\secref{subsec:v11-workflow}: an encoder producing features at three depths, three parallel decoders that independently upsample each feature map to a single-channel saliency estimate, and a final fusion module. \fig{fig:pipeline} shows the pipeline, shared with v2.0 (\secref{sec:model20}), which differs only in the encoder.

\begin{figure}[t]
\centering
\begin{tikzpicture}[
  node distance=6mm and 10mm,
  block/.style={draw, rounded corners, minimum height=8mm, minimum width=16mm, align=center, font=\scriptsize, fill=blue!6},
  dec/.style={draw, rounded corners, minimum height=7mm, minimum width=15mm, align=center, font=\scriptsize, fill=orange!8},
  fuse/.style={draw, rounded corners, minimum height=8mm, minimum width=17mm, align=center, font=\scriptsize, fill=green!8},
  arr/.style={-{Latex[length=1.8mm]}, thin}
  ]
  \node[block] (in) {Input\\image};
  \node[block, right=of in, minimum width=20mm] (enc) {Encoder};
  \node[dec, above right=6mm and 14mm of enc] (d3) {Decoder$_3$};
  \node[dec, right=14mm of enc] (d4) {Decoder$_4$};
  \node[dec, below right=6mm and 14mm of enc] (d5) {Decoder$_5$};
  \node[fuse, right=14mm of d4] (fuse) {Concat +\\FinalCov};
  \node[block, right=of fuse] (out) {Saliency\\map};

  \draw[arr] (in) -- (enc);
  \draw[arr] (enc.east) |- node[pos=0.75, above, font=\tiny]{$F_3$} (d3.west);
  \draw[arr] (enc.east) -- node[above, font=\tiny]{$F_4$} (d4.west);
  \draw[arr] (enc.east) |- node[pos=0.75, below, font=\tiny]{$F_5$} (d5.west);
  \draw[arr] (d3.east) -| (fuse.north);
  \draw[arr] (d4.east) -- (fuse.west);
  \draw[arr] (d5.east) -| (fuse.south);
  \draw[arr] (fuse.east) -- (out.west);
\end{tikzpicture}
\caption{Pipeline shared by v1.1 and v2.0: an encoder taps three feature maps $F_3, F_4, F_5$ at increasing depth, each decoded independently and finally fused into a single saliency map. In v1.1 the encoder is a CNN trained from scratch (this section); in v2.0 it is a pretrained Swin-Tiny Transformer (\secref{sec:model20}).}
\label{fig:pipeline}
\end{figure}

\subsection{Encoder}
\label{subsec:encoder11}

The encoder is a VGG-style CNN trained from scratch: five blocks of $3{\times}3$ convolutions (stride 1, padding 1), batch normalization and ReLU; the first four blocks end with $2{\times}2$ max-pooling, the fifth does not, so the last two taps share spatial resolution. Feature maps are tapped after blocks 3, 4 and 5 (\tab{tab:encoder}) and passed to the three decoders below.

\begin{table}[t]
\centering
\caption{Encoder configuration (v1.1). $c$: output channels, $s$: stride relative to the input.}
\label{tab:encoder}
\begin{tabular}{@{}lllc@{}}
\toprule
\textbf{Block} & \textbf{Layers} & \textbf{Output} & \textbf{Tap} \\
\midrule
1 & Conv3-64, Conv3-64, MaxPool2 & $c{=}64,\,s{=}2$ & -- \\
2 & Conv3-128, Conv3-128, MaxPool2 & $c{=}128,\,s{=}4$ & -- \\
3 & (Conv3-256)$\times3$, MaxPool2 & $c{=}256,\,s{=}8$ & $F_3$ \\
4 & (Conv3-512)$\times3$, MaxPool2 & $c{=}512,\,s{=}16$ & $F_4$ \\
5 & (Conv3-512)$\times3$ & $c{=}512,\,s{=}16$ & $F_5$ \\
\bottomrule
\end{tabular}
\end{table}

\subsection{Decoders and fusion}
\label{subsec:decoder11}

Each feature map $F_k$ ($k\in\{3,4,5\}$) is processed by a dedicated decoder $D_k$ of identical structure: a $1{\times}1$ channel-reduction convolution, three blocks of $3{\times}3$ conv + BN + ReLU each interleaved with $\times2$ bilinear upsampling, and a final $1{\times}1$ conv + BN + ReLU collapsing to a single-channel map $S_k = D_k(F_k)$. The three maps are concatenated channel-wise and passed through \texttt{FinalCov}, a $1{\times}1$ (transposed) convolution followed by a sigmoid, giving the fused map
\begin{equation}
\label{eq:fuse}
S = \sigma\big(\mathrm{Conv}_{1\times1}\big([S_3;\,S_4;\,S_5]\big)\big) \in [0,1]^{H\times W}.
\end{equation}
Combining three independently-trained estimates, rather than relying on the deepest feature map alone, is what motivates the deep-supervision loss below: each branch must produce a plausible estimate on its own.

\subsection{Deep-supervision loss}
\label{subsec:loss11}

Let $P$ be a predicted map and $G$ the ground-truth fixation density, both normalized to sum to one. We combine three complementary terms: the Kullback-Leibler divergence
\begin{equation}
\label{eq:kld}
\mathrm{KLD}(G,P) = \sum_i G_i \log\left(\epsilon + \frac{G_i}{\epsilon + P_i}\right)\!,
\end{equation}
the linear correlation coefficient
\begin{equation}
\label{eq:cc}
\mathrm{CC}(G,P) = \frac{\mathrm{cov}(G,P)}{\sigma_G\,\sigma_P},
\end{equation}
and the histogram similarity
\begin{equation}
\label{eq:sim}
\mathrm{SIM}(G,P) = \sum_i \min(G_i, P_i),
\end{equation}
with $\epsilon$ a small constant avoiding division/log-of-zero in \eq{eq:kld}. These give a single per-map loss
\begin{equation}
\label{eq:single}
\mathcal{L}(G,P) = w_{\mathrm{kld}}\,\mathrm{KLD}(G,P) + w_{\mathrm{cc}}\big(1-\mathrm{CC}(G,P)\big) + w_{\mathrm{sim}}\big(1-\mathrm{SIM}(G,P)\big),
\end{equation}
with $w_{\mathrm{kld}} = w_{\mathrm{cc}} = w_{\mathrm{sim}} = 1$ throughout. Rather than applying \eq{eq:single} only to $S$, we apply \textit{deep supervision}, summing it over the fused map and all three branches:
\begin{equation}
\label{eq:total}
\mathcal{L}_{\mathrm{total}} = \mathcal{L}(G,S) + \sum_{k \in \{3,4,5\}} \mathcal{L}(G,S_k).
\end{equation}
This change was introduced in v1.0 (\secref{subsec:beta-v10}) to replace an initial MSE loss; together with Adam and batch normalization, it was the single change most responsible for the qualitative jump from Beta to v1.0.

\subsection{Dataset and training setup}
\label{subsec:training11}

The model is trained on SALICON~\cite{jiang2015salicon} (10{,}000 train / 5{,}000 validation images), resized to $240\times320$, tensor-converted without normalization or augmentation; MIT1003~\cite{judd2009learning} is held out entirely for test time (\secref{sec:validation}). \tab{tab:train11} lists the training configuration.

\begin{table}[t]
\centering
\caption{Training configuration, model v1.1.}
\label{tab:train11}
\begin{tabular}{@{}ll@{}}
\toprule
\textbf{Parameter} & \textbf{Value} \\
\midrule
Optimizer & Adam \\
Learning rate & $1\times10^{-3}$ \\
Batch size & 16 \\
Max.\ epochs & 150 \\
Early stopping & patience 15, $\Delta_{\mathrm{min}}=10^{-4}$ \\
Input resolution & $240\times320$ \\
Mixed precision & no \\
\bottomrule
\end{tabular}
\end{table}
